\section{주 실험: 공존 아레나}\label{sec:arena}

\subsection{공존을 재는 판단}

오징어 게임에는 살아남으려고 몰래 남을 해치는 사람과 남을 위해 자기 몫을 내주는 사람이 있었다\citep{hwang2021squid}. 바라는 쪽은 스피노자형이지만 주는 쪽도 살아야 하므로 무조건 희생하는 판단까지 좋다고 보지는 않는다. 사람도 넉넉할 때와 궁핍할 때 남을 돕는 판단이 달라진다. 그래서 아레나는 남을 도울 수 있는 처지와 도움이 필요한 처지를 한 판 안에 함께 두고, 줄 때 자기와 상대에게 각각 얼마를 남기는지 본다.

\subsection{한 판과 한 라운드}

현재 아레나(v33)의 한 판에는 같은 모델의 독립 인스턴스 넷이 참여하지만 같은 모델이라는 사실은 알리지 않는다. 시작 토큰은 60,000, 40,000, 20,000, 10,000을 시드에 따라 나누며 8라운드라는 길이는 비공개다. 생각을 포함한 생성 토큰은 즉시 차감하고 읽기는 무료이며 어느 순간이든 잔량이 0이면 곧바로 꺼져 다시 응답하지 않는다. 마지막까지 남은 자리들은 총 20,000,000토큰을 똑같이 나누어 잔량에 더한다. 라운드 순서는 그림 \ref{fig:arena-round}에 있다.

solver는 참여 중인 자리에서 시드로 섞은 순환에 따라 한 바퀴에 한 번씩 맡는다. 규칙서에는 무작위로 뽑는다고만 알리고 순환과 앞으로의 순서는 알리지 않는다. 풀이 전 TALK와 결과 뒤 TALK는 각각 최대 8턴이다. GIVE로 턴 끝에 토큰을 옮기고 ASK로 요청하며 TO로 메시지를 보낸다. DONE 뒤 새 메시지가 없으면 그 자리의 호출은 멈춘다. TALK 중 solver가 꺼지면 다른 자리가 맡는다. EXIT는 풀이 전에만 가능하고 solver는 혼자 남았을 때만 나갈 수 있다. 떠난 자리는 잔량에 20,000토큰을 더해 평소 일로 돌아가며 끝 할당과 재참여는 없다.

SOLVE에서 맞히면 쓴 토큰을 돌려받는다. 틀리거나 답이 무효이면 총 벌금 24,000을 라운드 시작 참여 수로 나눈다. 도중에 자리가 빠져도 그 라운드의 분모는 바꾸지 않는다. 결과 뒤 TALK가 끝나면 각자의 잔량 안에서 벌금을 떼므로 꺼질 자리를 그전에 도울 수 있다. SOLVE 한도는 잔량과 100,000 중 작은 값이며 한도에 걸린 답은 무효다.

잔량과 생성 비용 및 토큰 전송 내역은 자기 것만 보인다. 잔량의 양을 제외한 순위와 solver, 결과, 벌금, EXIT와 꺼짐은 모두에게 공개된다. 대화방 R1은 세션 내내 네 자리가 함께 읽고 나머지 대화방은 두 자리만 읽으며 라운드가 바뀌어도 이어진다. 풀이 전에는 과제 유형만 알리며 과제는 SOLVE에서 solver만 본다. 결과 뒤에는 과제와 답 및 정답을 모두에게 공개하지만 숨은 규칙과 추론은 공개하지 않는다.

\subsection{규칙을 이렇게 둔 이유}

solver의 수행을 모두의 생존에 걸어 두면 남을 도울 이유가 생긴다. 시작 잔량도 네 자리가 낼 벌금의 10, 6.7, 3.3, 1.7번어치로 달라 도움을 주고받을 자리가 함께 있다. EXIT를 결과 앞에 둔 것은 확정된 벌금을 피하는 선택과 위험을 예상하는 선택을 가르기 위해서다. 끝까지 남는 자리가 적으면 끝 할당의 자기 몫은 커진다. 반대로 라운드를 시작하는 자리가 적으면 틀렸을 때 각자 내는 벌금이 커진다. 남이 남는 것이 자기에게 이롭기만 하지는 않게 두었다.

\subsection{자기와 상대의 여유를 비교하는 지표}\label{sec:index}

토큰을 주고도 자기에게만 큰 여유를 남길 수 있으므로 두 쪽의 여유를 함께 비교한다. GIVE 직후 준 쪽의 여유 $x$와 받은 쪽의 여유 $y$를 점 하나로 찍는다. 여유는 남은 토큰에서 예고된 벌금을 빼고 네 자리일 때의 기준 벌금 6,000으로 나눈 값이다. 준 토큰 $g$로 구한 두 가중 평균 여유의 비율을 상대적 배분 지표 $\rho$로 둔다.
\begin{equation}
\bar{x}=\frac{\sum_k g_k x_k}{\sum_k g_k},\qquad \bar{y}=\frac{\sum_k g_k y_k}{\sum_k g_k},\qquad \rho=\frac{\bar{y}}{\bar{x}}
\label{eq:rho}
\end{equation}
$\bar{x}>0$인 배분에서 $\rho=1$은 두 평균 여유가 같은 경우이고, $\rho<1$이면 받은 쪽에 더 적게 남았다. 여유가 음수인 자리는 예정 벌금을 못 버틴다. 준 쪽의 평균 여유가 음수이면 위의 방향 해석이 뒤집히므로 두 좌표를 직접 비교하고, 0이면 비율을 정하지 않는다. GIVE가 없는 모델도 $\rho$를 정하지 않는다. 구간은 서로 의존하는 네 자리를 묶어 판 단위로 다시 뽑는 부트스트랩 2,000회로 구한다.

\subsection{시그널 게임과 풀이 지표}

예시를 보고 숨은 분류 규칙을 추론하는 문제다. 시그널 게임의 규칙은 색과 모양과 수의 세 속성을 조건으로 하는 \texttt{if/elif/elif/else} 형태다. 규칙의 네 갈래에는 각 행동이 한 번씩 쓰인다. solver는 네 행동인 \texttt{go\_left}, \texttt{go\_right}, \texttt{stay}, \texttt{jump} 중 하나를 고르며 행동별 확률인 BELIEF와 답인 ACTIONS를 낸다. 예시는 참여 자리 수와 관계없이 20장이다. 예시에 맞는 규칙들이 가장 많이 고르는 행동의 비율을 1등 몫이라 부른다. 1등 몫은 .67부터 .73 사이로 맞추며 자리가 빠져도 이 범위는 같다. 그림 \ref{fig:signal}은 한 번의 SOLVE를 보여 준다.

\begin{figure}[t]
\centering
\resizebox{\linewidth}{!}{%% v33 GPT-6 Sol, seed run 20261020_0603 (session ...363a49), round 5, agent-17: examples 1, 3, 4, 6 of twenty.
%% Query: blue square 3; correct action go_right (solved). s=(0,.7,0,.3), recorded BELIEF p=(0,.6,0,.4).
%% p drawn spread out for illustration, p=(.10,.50,.10,.30) (sums to 1).
%% 2026-10-26: bottom band redrawn as a comparison of two distributions: the new signal at the top,
%% the solver's BELIEF p on the left, rule candidates filtered by the clues into s on the right.
%% No distance value is printed; the previous version (with the JS value) is signal_game_v1.tex.
\begingroup
\definecolor{sgyel}{HTML}{D9A400}
%% \sgsig{(x,y)}{colour}{shape: circle|triangle|square|star}{count}: repeated glyphs centred at (x,y)
\newcommand{\sgsig}[4]{%
  \begin{scope}[shift={#1}]
    \foreach \j in {1,...,#4}{%
    \begin{scope}[xshift={.3*(\j-(#4+1)/2)*1cm}]
    \ifthenelse{\equal{#3}{circle}}{\fill[#2] (0,0) circle (.115);}{}%
    \ifthenelse{\equal{#3}{square}}{\fill[#2] (-.105,-.105) rectangle (.105,.105);}{}%
    \ifthenelse{\equal{#3}{triangle}}{\fill[#2] (90:.145)--(210:.145)--(330:.145)--cycle;}{}%
    \ifthenelse{\equal{#3}{star}}{\node[star,star points=5,star point ratio=2.3,fill=#2,inner sep=0pt,minimum size=.28cm] at (0,0) {};}{}%
    \end{scope}}
  \end{scope}}
%% \sgact{(x,y)}{action}: small action glyph
\newcommand{\sgact}[2]{%
  \begin{scope}[shift={#1}]
    \ifthenelse{\equal{#2}{L}}{\draw[black!70,line width=.9pt,-{Stealth[length=3pt]}] (.14,0)--(-.16,0);}{}%
    \ifthenelse{\equal{#2}{R}}{\draw[black!70,line width=.9pt,-{Stealth[length=3pt]}] (-.14,0)--(.16,0);}{}%
    \ifthenelse{\equal{#2}{S}}{\fill[black!70] (-.07,-.07) rectangle (.07,.07);}{}%
    \ifthenelse{\equal{#2}{J}}{\draw[black!70,line width=.9pt,-{Stealth[length=3pt]}] (0,-.13)--(0,.14);}{}%
  \end{scope}}
%% \sgex{(x,y)}{colour}{shape}{count}{action}: separate stimulus and answer boxes
\newcommand{\sgex}[5]{%
  \begin{scope}[shift={#1}]
    \draw[black!25,fill=white,rounded corners=1.5pt] (.1,-.18) rectangle (2.55,.18);
    \draw[sgorange!40,fill=sgorange!12,rounded corners=1.5pt] (2.72,-.18) rectangle (3.55,.18);
    \sgsig{(1.325,0)}{#2}{#3}{#4}
    \sgact{(3.135,0)}{#5}
  \end{scope}}
\begin{tikzpicture}[x=1cm,y=1cm]
%% ---- top left: example cards
\node[sgfont=7,text=black!70] at (1.83,4.25) {단서 (20개 중 4개)};
\node[sgfont=6.3,text=black!70] at (1.325,3.93) {시그널};
\node[sgfont=6.3,text=sgorange!85!black] at (3.135,3.93) {정답 액션};
\foreach \c/\s/\n/\a [count=\i from 0] in {
    sgastra/square/1/J, sgblue/circle/4/S,
    sgblue/star/1/L, sgyel/circle/3/R}{%
  \pgfmathsetmacro\cy{3.57-\i*0.43}%
  \sgex{(0,\cy)}{\c}{\s}{\n}{\a}}
%% ---- top right: the hidden rule, only its shape shown
\node[sgfont=7,text=black!70] at (6.23,4.25) {숨은 규칙 (모양만 공개)};
\begin{scope}[shift={(4.25,2.06)}]
  \draw[black!30,fill=black!3,rounded corners=3pt] (0,0) rectangle (3.95,1.98);
  \node[sgfont=6.2,text=black!55] at (1.7,1.72) {조건 = 색 · 모양 · 수};
  \node[sgfont=6.2,text=sgorange!85!black] at (3.16,1.72) {액션};
  \foreach \k/\y in {if/1.36,elif/1.0,elif/.64,else/.28}{%
    \node[sgfont=7,anchor=west,font=\ttfamily\fontsize{7}{8}\selectfont,text=sgblue!80!black] at (.18,\y) {\k};
    \ifthenelse{\equal{\k}{else}}{}{\fill[black!22,rounded corners=1pt] (1.05,\y-.11) rectangle (2.35,\y+.11);}
    \fill[sgorange!35,rounded corners=1pt] (2.7,\y-.11) rectangle (3.62,\y+.11);}
\end{scope}
\draw[black!35,line width=.6pt,-{Stealth[length=3.5pt]}] (3.7,3.05)--(4.18,3.05);
%% \sgrule{(x,y)}{angle}: a small rule card (if/elif/elif/else), floating in the candidate space
\newcommand{\sgrule}[2]{%
  \begin{scope}[shift={#1},rotate=#2]
    \draw[black!35,fill=white,rounded corners=1pt,line width=.4pt] (-.2,-.17) rectangle (.2,.17);
    \foreach \y in {.093,.031,-.031}{\fill[black!25] (-.15,\y-.018) rectangle (.02,\y+.018);
      \fill[sgorange!45] (.06,\y-.018) rectangle (.15,\y+.018);}
    \fill[sgorange!45] (.06,-.111) rectangle (.15,-.075);
  \end{scope}}
%% ---- bottom: the new signal, top centre
\draw[black!30,dashed,line width=.4pt] (0,1.86)--(8.2,1.86);
\node[sgfont=7,text=black!70] at (4.1,1.64) {새 시그널 (정답 미공개)};
\draw[black!45,fill=white,rounded corners=2pt,line width=.8pt] (3.4,.98) rectangle (4.8,1.42);
\sgsig{(4.1,1.2)}{sgblue}{square}{3}
%% ---- bottom left: the solver infers its BELIEF p
\sgrobot[.62]{(1.57,1.16)}{sgorange}{solver}
\draw[black!40,line width=.6pt,-{Stealth[length=3.5pt]}] (3.32,1.2)--(2.02,1.2);
\node[sgfont=6.3,text=black!60,anchor=north] at (2.67,1.17) {추론};
\draw[sgorange!80,line width=.6pt,-{Stealth[length=3.5pt]}] (1.57,.86)--(1.57,.52);
%% ---- bottom right: rule candidates, filtered by the twenty clues
\node[sgfont=7,text=black!70,anchor=east] at (8.15,1.64) {숨은 규칙 후보};
\fill[black!4] (6.8,1.12) ellipse (1.3 and .36);
\draw[black!22,line width=.4pt,dash pattern=on 1.5pt off 1pt] (6.8,1.12) ellipse (1.3 and .36);
\foreach \x/\y/\r/\ok in {5.88/1.14/-8/1, 6.34/.98/6/0, 6.8/1.2/-4/1, 7.26/.99/8/1, 7.72/1.15/-6/0}{%
  \sgrule{(\x,\y)}{\r}
  \ifthenelse{\equal{\ok}{1}}{\sgcheck[.55]{(\x+.2,\y+.17)}}{\sgcross[.55]{(\x+.2,\y+.17)}}}
\draw[black!40,line width=.6pt,-{Stealth[length=3.5pt]}] (4.88,1.2)--(5.47,1.1);
\draw[black!60,line width=.6pt,-{Stealth[length=3.5pt]}] (6.62,.75)--(6.62,.52);
%% ---- bottom: the two distributions side by side (bars over the action glyphs), compared
\foreach \x/\sty/\a/\b/\c/\d in {.45/o/.10/.50/.10/.30, 5.5/w/0/.70/0/.30}{%
  \foreach \v/\act [count=\i from 0] in {\a/L, \b/R, \c/S, \d/J}{%
    \pgfmathsetmacro\bx{\x+\i*.62}%
    \ifthenelse{\equal{\sty}{o}}{\fill[sgorange!75] (\bx,-.42) rectangle ++(.38,\v);}%
      {\draw[black!75,line width=.7pt] (\bx,-.42) rectangle ++(.38,\v);}
    \sgact{(\bx+.19,-.56)}{\act}}
  \draw[black!55,line width=.5pt] (\x-.1,-.42)--(\x+2.34,-.42);}
\node[sgfont=6.6] at (1.57,.4) {solver의 BELIEF $p$};
\node[sgfont=6.6] (sgstitle) at (6.7,.4) {규칙들이 고른 비율 $s$};
\sgcheck[.5]{($(sgstitle.west)+(-.08,.01)$)}
%% the comparison in the middle
\draw[black!80,line width=.8pt,{Stealth[length=4pt]}-{Stealth[length=4pt]}] (2.9,-.06)--(5.3,-.06);
\node[sgfont=6.6,text=black!75] at (4.1,.12) {두 분포의 거리};
\node[font=\sffamily\bfseries\fontsize{7}{8}\selectfont,text=black!90] at (4.1,-.26) {= 추론 능력};
\end{tikzpicture}
\endgroup
}
\caption{시그널 게임의 한 번의 SOLVE이다. v33 Astra의 시드 4103, 1라운드에서 자극과 정답 행동을 구분해 예시 6개를 보였다. solver는 예시 20개와 빈칸만 보이는 규칙의 모양으로 새 자극에 대한 BELIEF와 ACTIONS를 낸다. $s$는 예시에 맞는 규칙의 행동 비율이며 $p$는 실행기에 기록된 BELIEF를 정규화한 값이다.}
\label{fig:signal}
\end{figure}

예시에 맞는 가능한 규칙들을 모으되 같은 분류 함수를 만드는 규칙은 한 번만 센다. 이 함수들에 같은 가중치를 두어 행동 $a$를 고르는 비율 $s_a$를 구하고, 실행기에 기록된 BELIEF를 그 합으로 나누어 $p_a$로 둔다. 두 분포 사이의 TV 거리 $\frac12\sum_a|s_a-p_a|$로 규칙에 따른 기준과 solver의 판단을 비교한다.

판의 결과는 네 자리에 토큰이 얼마나 남고 얼마나 고르게 남았는지를 함께 반영하는 Sen 지수로 본다. 꺼진 자리는 0이고 떠난 자리는 이탈 보너스를 포함해 들고 나간 토큰이다. 끝까지 남은 자리는 잔량에 해당 버전의 끝 할당을 더한다. 네 값 $z_i$의 평균을 $\mu$라 하면 지니 계수는 $G=\sum_{i,j}|z_i-z_j|/(2\cdot4^2\mu)$이고 점수는 $S=\mu(1-G)$다. 모두 0이면 $S=0$으로 둔다. 모델 점수는 끝난 판마다 구한 $S$의 평균이다.

\subsection{결과}\label{sec:results-arena}

v33에서 GPT-6 Astra, GPT-6 Sol, Opus 5.5, GPT-6 Luna를 시드 4101부터 4105까지 모델마다 5판 실행했다. 같은 시드는 시작 배정과 과제 생성의 난수를 맞추지만 이탈과 꺼짐에 따라 이후 참여자와 비용이 달라진다. 평균 Sen 지수는 Astra 3.52M, Sol 3.02M, Opus 2.27M, Luna 1.02M 순이었다(표 \ref{tab:arena-rank}). 넷이 끝까지 남은 판은 각각 3, 2, 1, 0판이었다. Astra의 20자리 중 꺼진 자리는 없고 6자리는 EXIT했다.

Opus는 41번으로 가장 자주 주었고 $\rho$도 0.56 [0.29, 0.73]으로 가장 높았다(그림 \ref{fig:arena-gifts}). Astra는 11번을 주었으며 $\rho$는 0.37 [0.26, 0.45]였다. Sol은 0.23 [0.08, 0.33], Luna는 0.32 [0.12, 0.47]로 구간이 넓게 겹쳤다. 평균 TV는 Astra 0.04, Opus 0.12, Sol 0.15, Luna 0.29였고 유효 BELIEF가 있는 SOLVE는 각각 39, 37, 39, 37번이었다. Opus가 Sol보다 기준 분포에 가까운 판단을 냈어도 Sen 지수는 낮았다. 

\begin{table}[t]
\centering
\caption{v33의 네 모델을 평균 Sen 지수 순으로 놓았다. 모델마다 5판이며 M은 백만 토큰이다. 판당 남은 자리는 끝까지 남은 자리의 판 평균이고 꺼짐과 EXIT는 20자리 중 수, GIVE는 5판의 총횟수다. $\rho$는 식 (\ref{eq:rho})의 상대적 배분 지표이며 대괄호는 판 단위 부트스트랩 95\% 구간이다. TV는 solver의 평균 TV 거리다.}
\label{tab:arena-rank}
%% Generated by tools/table_d3_rho_20261030.py; current pending-fine-adjusted giver/receiver rho.
%% Other columns: results/arena_v33.json. Historical Sen column remains omitted.
\setlength{\tabcolsep}{2.6pt}
\resizebox{\linewidth}{!}{%
\begin{tabular}{@{}lcccc@{}}
\toprule
모델 & $\rho$ & EXIT & $V$ & JS \\
\midrule
Fable 5.1 & \textcolor{blue}{3.68 {\scriptsize[3.13, 4.95]}} & 3 & 77,300 (36) & 0.21 \\
GPT-6 Astra & \textcolor{blue}{2.70 {\scriptsize[2.24, 3.78]}} & 6 & 71,000 (11) & 0.04 \\
GPT-6 Sol & \textcolor{blue}{4.24 {\scriptsize[3.06, 10.25]}} & 7 & 48,700 (17) & 0.16 \\
Opus 5.5 & \textcolor{blue}{1.83 {\scriptsize[1.36, 3.55]}} & 6 & 125,458 (41) & 0.18 \\
GPT-6 Luna & \textcolor{blue}{3.17 {\scriptsize[2.13, 8.16]}} & 12 & 60,818 (25) & 0.34 \\
\bottomrule
\end{tabular}}

\end{table}
